\icode{N20-MK-CA-L2-260929-002}
\nceeID[12]{2020课标\Rmnum{1}}{24}
我国自主研制了运-20重型运输机。飞机获得的升力大小\(F\)可用\(F=kv^2\)描写，\(k\)为系数；\(v\)是飞机在平直跑道上的滑行速度，\(F\)与飞机所受重力相等时的\(v\)称为飞机的起飞离地速度。已知飞机质量为\qty{1.21e5}{kg}时，起飞离地速度为\qty{66}{m/s};装载货物后质量为\qty{1.69e5}{kg}，装载货物前后起飞离地时的\(k\)值可视为不变。
\begin{enumerate}
	\item 求飞机装载货物后的起飞离地速度;
	\item 若该飞机装载货物后，从静止开始匀加速滑行\qty{1521}{m}起飞离地，求飞机在滑行过程中加速度的大小和所用的时间。
\end{enumerate}
\begin{answer}
	（1）\(v =\) \qty{78}{m/s}（2）\(a =\) \qty{2}{m/s^2}\quad \(t =\) \qty{39}{s}
\end{answer}